\documentclass[10pt,a4paper]{amsart}
\usepackage[margin=19mm]{geometry}
\usepackage{amsmath,amssymb,mathtools}
\usepackage[scr=rsfs,cal=euler]{mathalfa}
\usepackage{xfrac}
\usepackage[unicode,hidelinks,bookmarks=false]{hyperref}
\newcommand{\ie}{i.e.\ }
\newcommand{\cf}{cf.\ }
\newcommand{\resp}{respectively\ }
\newcommand{\Subsection}{\subsection}
\providecommand{\nobreakdash}{\nobreak\hbox{-}}
\setlength{\emergencystretch}{2em}
\multlinegap0pt
\usepackage{tikz}
\usetikzlibrary{cd}
\usepackage{amssymb}
\newtheorem{thm}{Theorem}
\newcommand{\alg}{\mathrm{alg}}
\newcommand{\s}{\mathbb{S}}
\title{Homotopy properties of regular mappings into real retract rational varieties}
\author{Juliusz Banecki}
\date{Source: arXiv:2505.11203v3 (CC BY 4.0)}
\begin{document}
\maketitle

\noindent\emph{Source and licence.} Homotopy properties of regular mappings into real retract rational varieties. Version arXiv:2505.11203v3 (CC BY 4.0), distributed by arXiv under the Creative Commons Attribution 4.0 International licence, http://creativecommons.org/licenses/by/4.0/, as recorded in the arXiv licence metadata for this exact version and shown on the abstract page https://arxiv.org/abs/2505.11203v3. Full licence notice: \texttt{licenses/arxiv-2505.11203-CC-BY-4.0.txt}.\par\smallskip

\noindent\textit{Editorial scope.} Counted unit: the article's thm thm\_products (source0.tex lines 156--158). The article's complete original proof of it is delivered with it (source0.tex lines 301--335, one \texttt{proof} environment of 35 lines, which the article places away from the statement). No mathematics is written, repaired or re-derived: the statement and the proof are the article's own lines, in the article's own notation, and no retained line is substituted or re-flowed.  Retained uncounted as the source-local context the proof needs: lines 120--127.  Nothing was omitted from any retained span, so the package carries no omission record.  No retained line contains a citation, so the wrapper carries no bibliography and no bibliography entry.  The retained lines contain the article's own diagram code, which the wrapper typesets with the standard tikz packages; no image file is delivered and the package declares no asset dependency.  Editorial formatting only, in addition: an AMS \texttt{amsart} preamble that restates the article's own declarations and macros used by the retained lines, namely the environments \texttt{thm} on separate counters, together with the standard packages those lines need.  The retained environments print in this document as Theorem 1 in order of appearance, whereas the article prints the same items with the numbering of its own full document.  The complete native source of the article as distributed in the arXiv e-print for this exact version is bundled unchanged as \texttt{sources/178223/source0.tex}.
\begin{thm}\label{approx_thm}
Let $Y$ be a nonsingular retract rational affine variety, let $X$ be an affine variety, let $Z\subset X$ be a Zariski closed subvariety of $X$ compact in the Euclidean topology, and let $f:X\rightarrow Y$ be a continuous mapping satisfying the following two conditions:
\begin{enumerate}
    \item $f$ is homotopic to a regular mapping $g:X\rightarrow Y$,
    \item $f\vert_Z:Z\rightarrow Y$ is regular. 
\end{enumerate}
Then $f$ can be approximated by regular mappings $\widetilde{f}:X\rightarrow Y$ in the compact-open topology, agreeing with $f$ on $Z$.
\end{thm}

\begin{thm}\label{thm_products}
Let $Y$ be a nonsingular retract rational affine variety with a distinguished point $y_0\in Y$. Let $n,m\geq 1$ be natural numbers. Then, for any two classes $\alpha\in \pi_n(Y,y_0),\;\beta\in \pi_m(Y,y_0)$, their Whitehead product $[\alpha,\beta]\in\pi_{n+m-1}(Y,y_0)$ belongs to $\pi_{n+m-1}^\alg(Y,y_0)$.
\end{thm}

\begin{proof}[Proof of Theorem \ref{thm_products}]
By the definition, given two continuous mappings $f:(\s^n,e)\rightarrow (Y,y_0),g:(\s^m,e)\rightarrow (Y,y_0)$ the Whitehead product of their classes is represented by the following composition
\begin{center}
    \begin{tikzcd}
h:\s^{n+m-1} \arrow[r, "\phi"] & \s^n\vee\s^m \arrow[r, "f\vee g"] & Y
\end{tikzcd}
\end{center}
where $\phi$ is the attaching map of the $n+m$-cell in the standard CW-structure of the torus $\s^n\times \s^m$. Let $\text{const}_n:\s^n\rightarrow Y,\;\text{const}_m:\s^m\rightarrow Y$ denote the mappings constantly equal to $y_0$, and define:
\begin{align*}
    \psi:=(f\vee \text{const}_m)\circ \phi, \\
    \xi:=(\text{const}_n\vee g)\circ \phi.
\end{align*}

It is then clear that $\psi$ and $\xi$ are both homotopic to the constant mapping, as they represent the homotopy classes of the products $[f,\text{const}_m]$ and $[\text{const}_n, g]$ respectively. Moreover it follows from the construction that for every $x\in \s^n$ we have
\begin{equation*}
\text{either }
\begin{cases}
\psi(x)=h(x),\\
\xi(x)=y_0,
\end{cases}
\text{or }
\begin{cases}
\psi(x)=y_0,\\
\xi(x)=h(x).
\end{cases}
\end{equation*}
This shows that $h$ is homotopic to $\mu(\psi,\xi)$ relative to $\{e\}$.

Applying Theorem \ref{approx_thm} we can find regular mappings 
\begin{align*}
    \tilde{\psi}:\s^{n+m-1}\rightarrow Y,\\
    \tilde{\xi}:\s^{n+m-1}\rightarrow Y,
\end{align*}
which satisfy $\psi(e)=\xi(e)=y_0$ and approximate $\psi$ (respectively $\xi$) closely. If the approximation is sufficiently close then $\tilde{h}:=\mu(\tilde{\psi},\tilde{\xi})$ is a close regular approximation of $\mu(\psi,\xi)$. In particular it represents the same homotopy class.
\end{proof}

\end{document}
